\documentclass[11pt]{article}
\usepackage[margin=1in]{geometry}
\usepackage{amsmath,amssymb}
\usepackage{booktabs}
\usepackage{xcolor}
\usepackage[hidelinks]{hyperref}

\definecolor{hi}{RGB}{0,120,0}
\definecolor{lo}{RGB}{170,0,0}
\definecolor{hdr}{RGB}{30,30,90}
\newcommand{\code}[1]{\texttt{\small #1}}
\newcommand{\callout}[1]{%
  \begin{center}\fcolorbox{hdr}{blue!4}{%
  \begin{minipage}{0.95\textwidth}\vspace{2pt}#1\vspace{2pt}\end{minipage}}\end{center}}

\title{\vspace{-1.5em}\textbf{Thread 3 --- Chromatogram Integration: allocation analysis}\\[0.3em]
\large Why it is not a quick allocation win, unlike Quant \& Output}
\author{}
\date{2026-06-24 \quad\textbar\quad branch \code{perf/runtime-profiling}}

\begin{document}
\maketitle
\vspace{-2.5em}

\callout{
\textbf{Bottom line.} The Chromatogram Integration phase ($\sim$3.4\,GB cumulative
allocation, warm 2-file MTAC) is dominated by the \emph{actual chromatogram signal}
being materialized, plus one structure-of-arrays conversion copy. Unlike Quant \&
Output --- whose 3.9\,GB win came from skipping \emph{redundant} string churn (99.5\%
duplicate accessions) --- Integration's allocation is mostly \emph{necessary numeric
data} that cannot be skipped. The single reducible item (the conversion copy) needs a
moderate refactor, so it does \emph{not} clear the ``no excessive complexity'' bar as a
quick win. Recommendation: measure the copy precisely before deciding.
}

\section{What the phase does}
Three stages, in order:
\begin{enumerate}
  \item \textbf{\code{build\_chromatograms}} (\code{utils.jl:707}) --- walks the scan
  range, runs the fused deconvolution kernel per scan, and \textbf{emits one
  \code{MS2ChromObject} per (precursor, scan) chromatogram point} into a
  \code{Vector\{MS2ChromObject\}} (pre-sized to \code{n\_scans}$\times$100). It returns
  \code{DataFrame(@view(chromatograms[1:rt\_idx]))} (\code{utils.jl:878}).
  \item \textbf{\code{integrate\_precursors}} (\code{utils.jl:399}) --- groups the
  chromatogram table by precursor and integrates each peak.
  \item \textbf{\code{integrate\_chrom}} (\code{integrate\_chrom.jl}) --- per precursor:
  Whittaker--Henderson smoothing $\rightarrow$ second-derivative bounds $\rightarrow$
  baseline-subtracted trapezoidal area.
\end{enumerate}

\section{Where the 3.4\,GB is --- and is not}
\textbf{The integration loop is already allocation-lean.} \code{integrate\_precursors}
(stages 2--3) was written to reuse memory:
\begin{itemize}
  \item per-thread \code{WHWorkspace} and \code{Chromatogram} state are pre-allocated
  once (\code{utils.jl:431--432}) and reused across all precursors;
  \item every chromatogram slice is passed as an \code{@view} (\code{utils.jl:464--467})
  --- no per-precursor copies;
  \item results are written into caller-provided output vectors (\code{peak\_area},
  \code{new\_best\_scan}, \code{points\_integrated}).
\end{itemize}
So the per-precursor smoothing/integration is \emph{not} the allocator. (The line
profiler agreed: \code{WHSmooth!}~$\approx$0.11\,GB, \code{build\_chromatograms}
leaf~$\approx$0.15\,GB --- small.)

\textbf{The allocation lives in stage 1, and it is mostly real data:}
\begin{itemize}
  \item \textcolor{lo}{\textbf{the \code{Vector\{MS2ChromObject\}} of chromatogram
  points}} --- this \emph{is} the signal being integrated (every (precursor, scan)
  intensity in each precursor's RT window). It scales with the data, not with waste.
  \item \textcolor{lo}{\textbf{the AoS$\rightarrow$DataFrame conversion copy}}
  (\code{DataFrame(@view(...))}, \code{utils.jl:878}) --- the array-of-structs
  \code{MS2ChromObject} vector is copied field-by-field into the \code{DataFrame}'s
  columns; the original vector is then discarded. This is a transient $\sim$1$\times$
  duplicate of the chromatogram data, analogous to the Main Search
  \code{df\_materialize} (0.78\,GB there).
\end{itemize}

\section{Why this is unlike the Quant \& Output wins}
\begin{center}
\begin{tabular}{lll}
\toprule
 & \textbf{Quant \& Output (Thread 2)} & \textbf{Chrom Integration (Thread 3)} \\
\midrule
What allocates & per-row \code{split} + \code{String} interning & chromatogram signal + AoS$\to$DF copy \\
Nature & \textcolor{lo}{redundant} (99.5\% duplicate rows) & \textcolor{hi}{necessary} (numeric data) \\
Fix & memoize \& skip repeats & emit columns (SoA) instead of copy \\
Complexity & low (a \code{Set}/cache) & moderate (restructure emit path) \\
Result & \textbf{$-$92\% / $-$55\%, bit-identical} & only the copy is reducible \\
\bottomrule
\end{tabular}
\end{center}
\noindent Thread 2 worked because the allocation was \emph{repeated work over duplicate
strings} --- skippable with a tiny memo. Thread 3 has no such redundancy: the
chromatogram points are distinct measured intensities. The only thing we are paying
\emph{twice} for is the AoS vector plus its DataFrame copy.

\section{The one available lever, and the recommendation}
Avoiding the conversion copy means having \code{build\_chromatograms} write directly
into column vectors (\code{rt}, \code{scan\_idx}, \code{intensity},
\code{precursor\_fraction\_transmitted}, \ldots) --- a structure-of-arrays emit path ---
and constructing the \code{DataFrame} from those columns with no per-element copy. That
removes the duplicate but is a \textbf{moderate refactor} of the emit hot loop (and must
stay bit-identical and not regress speed).

\callout{
\textbf{Recommendation.} Before refactoring, \emph{measure}: add \code{gc\_bytes}
sub-attribution (same method as the Main Search \code{MSLOOP} pass) to split the 3.4\,GB
into (a) the chromatogram \code{Vector}, (b) the DataFrame conversion copy, and (c)
everything else. Refactor \emph{only} if the copy is a clean $\gtrsim$1\,GB and the SoA
emit stays simple. Otherwise, conclude Thread 3: the phase is near its necessary-data
floor and further cuts would trade real complexity for little memory.
}

\section{Sub-attribution results (measured) --- the hypothesis was wrong}
\code{gc\_bytes} bucketing (env \code{PIONEER\_DIAG\_CHROM}, single-thread, 1-file MTAC,
2{,}726{,}965 chromatogram rows):
\begin{center}
\begin{tabular}{lrl}
\toprule
\textbf{Step} & \textbf{GB} & \textbf{Note} \\
\midrule
\code{integrate\_precursors} & 0.724 & \textbf{largest} --- inside \code{integrate\_chrom} (per precursor) \\
\code{extract\_chromatograms} (Vector build) & 0.592 & per-scan deconv + the \code{MS2ChromObject} signal \\
\code{sort\_chromatograms\_for\_integration!} & 0.189 & \code{sortperm} \\
\code{df\_convert} (AoS$\rightarrow$DataFrame copy) & \textcolor{lo}{\textbf{0.074}} & \textcolor{lo}{\textbf{the hypothesized target --- negligible}} \\
\code{get\_isotopes\_captured!} & 0.018 & small \\
\bottomrule
\end{tabular}
\end{center}
\callout{
\textbf{Corrected conclusion.} The AoS$\rightarrow$DataFrame copy is only \textbf{74\,MB} ---
the measure-first guard worked, the SoA refactor is \emph{not} worth it. The real targets
are \code{integrate\_chrom} (0.72\,GB, $\sim$2.4\,KB/precursor across $\sim$300k passing
PSMs; it defines 9 nested helper functions --- a possible per-call closure-capture cost,
same pattern the team fixed in \code{process\_scans\_fused}) and the \code{extract}
Vector build (0.59\,GB, partly the necessary signal). Next step: drill into
\code{integrate\_chrom} the same way before any fix --- do \emph{not} assume the closures
are the cause until measured.
}

\end{document}
